\documentclass[letterpaper,11pt]{article}
\usepackage[top=1in,bottom=1in,left=1in,right=1in]{geometry}
\usepackage[T1]{fontenc}
\usepackage[utf8]{inputenc}
\usepackage[spanish,es-nodecimaldot]{babel}
\usepackage{lmodern,microtype,xcolor,tcolorbox,tabularx,booktabs}
\usepackage[hidelinks]{hyperref}

\definecolor{colorHanniel}{HTML}{5941A9}
\definecolor{colorAlan}{HTML}{166A8F}
\definecolor{colorZoe}{HTML}{A52F64}
\definecolor{colorLuis}{HTML}{9A4C16}
\definecolor{colorJoseph}{HTML}{197349}
\definecolor{colorAlex}{HTML}{67531C}

\title{\textbf{Guion Técnico Audiovisual — Práctica 6}\\\large Empaquetado JAR, Javadoc y Modularización del Curso de POO}
\author{\textbf{Equipo 2} — Grupo 7, Facultad de Ingeniería, UNAM}
\date{Octubre 2026}

\begin{document}
\maketitle

\section*{Ficha Técnica de Producción}
\begin{itemize}
    \item \textbf{Duración Total del Video:} 05:00 minutos (300 segundos exactos).
    \item \textbf{Introducción Externa Institucional:} 00:00 -- 00:42 (42 segundos, sin diálogo redactado).
    \item \textbf{Diálogo Continuo:} 00:42 -- 04:42 (240 segundos, 457 palabras en 11 bloques alternados).
    \item \textbf{Cierre Visual:} 04:42 -- 05:00 (18 segundos, sin diálogo).
    \item \textbf{Elenco y Roles:} Hanniel (Cubo de Geometry Dash, 80 s en 6 bloques), Alan (CJ de GTA, 32 s), Zoé (Snoopy, 32 s), Luis (Rabbit, 32 s), Joseph (Rick Sanchez, 32 s), Alex (Phineas, 32 s).
\end{itemize}

\vspace{0.5em}
\noindent\hrulefill
\vspace{0.5em}

\subsection*{Bloque 1 — 00:42 a 00:56 (14 s) · Referencia: \texttt{1.mp3}}
\textbf{Hablante:} Hanniel \textcolor{colorHanniel}{(Cubo de Geometry Dash, \#5941A9)} \hfill \textbf{Tono:} Entusiasta y acogedor\\
\textbf{Visual / Transición:} Transición de cortinilla suave; caja digital con el logo de Java cerrándose.\\
\textbf{Sustento Técnico:} Sección 1 (Introducción y planteamiento del empaquetado modular).
\begin{tcolorbox}[colback=colorHanniel!5,colframe=colorHanniel,title=\textbf{Locución Definitiva — Hanniel}]
¡Hola a todos! Bienvenidos. Hoy les contamos cómo ordenamos todos nuestros programas de Java del semestre usando empaquetado JAR y documentación con Javadoc. Alan, ¿por qué era indispensable este cambio?
\end{tcolorbox}

\subsection*{Bloque 2 — 00:56 a 01:28 (32 s) · Referencia: \texttt{2.mp3}}
\textbf{Hablante:} Alan \textcolor{colorAlan}{(CJ de GTA, \#166A8F)} \hfill \textbf{Tono:} Claro, explicativo y reflexivo\\
\textbf{Visual / Transición:} Corte directo; diagrama de múltiples clases compiladas dispersas agrupándose dentro de un archivo \texttt{.jar}.\\
\textbf{Sustento Técnico:} Sección 2.1 (Concepto, estructura y distribución de archivos JAR).
\begin{tcolorbox}[colback=colorAlan!5,colframe=colorAlan,title=\textbf{Locución Definitiva — Alan}]
Antes teníamos decenas de archivos sueltos por todas partes. Si querías compartir un ejercicio o probarlo en otra máquina, era fácil olvidar una clase o perder la estructura. El archivo JAR resuelve esto empaquetando todo el proyecto en un único contenedor listo para ejecutarse, simplificando la distribución y garantizando que ningún archivo esencial quede fuera del sistema final.
\end{tcolorbox}

\subsection*{Bloque 3 — 01:28 a 01:42 (14 s) · Referencia: \texttt{3.mp3}}
\textbf{Hablante:} Hanniel \textcolor{colorHanniel}{(Cubo de Geometry Dash, \#5941A9)} \hfill \textbf{Tono:} Dinámico y conectivo\\
\textbf{Visual / Transición:} Paneó dinámico hacia visualización de árbol de directorios con nombres de paquetes formales.\\
\textbf{Sustento Técnico:} Sección 3.2 (Organización modular bajo \texttt{mx.unam.fi.die.poo.g7}).
\begin{tcolorbox}[colback=colorHanniel!5,colframe=colorHanniel,title=\textbf{Locución Definitiva — Hanniel}]
Exacto, Alan. Un JAR es como una caja de envíos sellada: reúne todo el código para que viaje y funcione sin desarmarse. Para mantener el orden interno, Zoé, ¿cómo organizamos los paquetes?
\end{tcolorbox}

\subsection*{Bloque 4 — 01:42 a 02:14 (32 s) · Referencia: \texttt{4.mp3}}
\textbf{Hablante:} Zoé \textcolor{colorZoe}{(Snoopy, \#A52F64)} \hfill \textbf{Tono:} Didáctico, amable y seguro\\
\textbf{Visual / Transición:} Gráfico ilustrado de casilleros identificados con p1 hasta p5 y proyecto, resaltando el espacio de nombres institucional.\\
\textbf{Sustento Técnico:} Sección 2.3 y 3.2 (Convención institucional de nombres y jerarquía de paquetes).
\begin{tcolorbox}[colback=colorZoe!5,colframe=colorZoe,title=\textbf{Locución Definitiva — Zoé}]
Para que no se mezclen los programas, creamos paquetes con el formato oficial de la facultad. Es como asignarle a cada práctica un casillero único: p1 para la calculadora, hasta p5 para empleados y el proyecto del hospital. Así evitamos nombres repetidos y logramos que cada módulo viva en su propio espacio bien identificado y protegido.
\end{tcolorbox}

\subsection*{Bloque 5 — 02:14 a 02:27 (13 s) · Referencia: \texttt{5.mp3}}
\textbf{Hablante:} Hanniel \textcolor{colorHanniel}{(Cubo de Geometry Dash, \#5941A9)} \hfill \textbf{Tono:} Interesado y motivador\\
\textbf{Visual / Transición:} Zoom digital hacia un libro de especificaciones abriéndose sobre una pantalla de computadora.\\
\textbf{Sustento Técnico:} Sección 2.2 (Propósito de la documentación con Javadoc).
\begin{tcolorbox}[colback=colorHanniel!5,colframe=colorHanniel,title=\textbf{Locución Definitiva — Hanniel}]
Así es, Zoé. Cada práctica tiene su casillero institucional. Pero además del orden, necesitábamos que cualquiera entienda el funcionamiento de cada pieza. Luis, ¿cómo nos ayuda Javadoc en esta tarea?
\end{tcolorbox}

\subsection*{Bloque 6 — 02:27 a 02:59 (32 s) · Referencia: \texttt{6.mp3}}
\textbf{Hablante:} Luis \textcolor{colorLuis}{(Rabbit, \#9A4C16)} \hfill \textbf{Tono:} Amigable, analítico y pedagógico\\
\textbf{Visual / Transición:} Transición lateral; gráfico analógico contrastando un aparato sin manual frente a una guía técnica estructurada.\\
\textbf{Sustento Técnico:} Sección 2.2 (Documentación técnica y lenguaje natural con etiquetas Javadoc).
\begin{tcolorbox}[colback=colorLuis!5,colframe=colorLuis,title=\textbf{Locución Definitiva — Luis}]
Imagina que compras un electrodoméstico sin manual de instrucciones. Sería un desastre entenderlo. Javadoc es justamente ese manual técnico para programadores: explica en lenguaje natural qué hace cada función, qué datos necesita y qué resultados produce. Al estandarizar los comentarios, transformamos líneas de código abstracto en una guía de consulta clara y navegable para cualquier desarrollador.
\end{tcolorbox}

\subsection*{Bloque 7 — 02:59 a 03:12 (13 s) · Referencia: \texttt{7.mp3}}
\textbf{Hablante:} Hanniel \textcolor{colorHanniel}{(Cubo de Geometry Dash, \#5941A9)} \hfill \textbf{Tono:} Enfocado y entusiasta\\
\textbf{Visual / Transición:} Animación de etiquetas doclet transformándose automáticamente en páginas web HTML locales.\\
\textbf{Sustento Técnico:} Sección 3.1 (Investigación técnica sobre la generación óptima de Javadoc).
\begin{tcolorbox}[colback=colorHanniel!5,colframe=colorHanniel,title=\textbf{Locución Definitiva — Hanniel}]
Gran analogía, Luis. Javadoc convierte los comentarios estructurados en un manual web interactivo para desarrolladores. Joseph, tú que analizaste las herramientas, ¿cuál es la mejor forma técnica de generarlo?
\end{tcolorbox}

\subsection*{Bloque 8 — 03:12 a 03:44 (32 s) · Referencia: \texttt{8.mp3}}
\textbf{Hablante:} Joseph \textcolor{colorJoseph}{(Rick Sanchez, \#197349)} \hfill \textbf{Tono:} Técnico, directo y pragmático\\
\textbf{Visual / Transición:} Gráfico de consola ejecutando la herramienta javadoc con UTF-8 y generando carpetas individuales sin errores.\\
\textbf{Sustento Técnico:} Sección 3.1 y 4 (Generación determinista en consola, classpath y sourcepath).
\begin{tcolorbox}[colback=colorJoseph!5,colframe=colorJoseph,title=\textbf{Locución Definitiva — Joseph}]
La mejor manera es invocar la herramienta oficial javadoc desde la consola de comandos, especificando codificación universal UTF-8 y aislando cada clase. Aunque los entornos visuales ofrecen botones rápidos, la terminal garantiza un control estricto de rutas y dependencias. De este modo, generamos documentación determinista y libre de advertencias en las carpetas destinadas para cada componente.
\end{tcolorbox}

\subsection*{Bloque 9 — 03:44 a 03:57 (13 s) · Referencia: \texttt{9.mp3}}
\textbf{Hablante:} Hanniel \textcolor{colorHanniel}{(Cubo de Geometry Dash, \#5941A9)} \hfill \textbf{Tono:} Conector y propositivo\\
\textbf{Visual / Transición:} Transición circular; diagrama global conectando las cinco prácticas hacia el núcleo del proyecto de hospital.\\
\textbf{Sustento Técnico:} Sección 3.3 (Integración de modularidad y avance del proyecto hospitalario).
\begin{tcolorbox}[colback=colorHanniel!5,colframe=colorHanniel,title=\textbf{Locución Definitiva — Hanniel}]
Totalmente, Joseph. La consola garantiza control absoluto de codificación y dependencias. Y para concluir este recorrido, Alex, ¿cómo se conecta todo este trabajo con nuestro proyecto final?
\end{tcolorbox}

\subsection*{Bloque 10 — 03:57 a 04:29 (32 s) · Referencia: \texttt{10.mp3}}
\textbf{Hablante:} Alex \textcolor{colorAlex}{(Phineas, \#67531C)} \hfill \textbf{Tono:} Concluyente, optimista y profesional\\
\textbf{Visual / Transición:} Diagrama del sistema hospitalario integrando médicos, enfermeros y pacientes bajo su paquete y ejecutable propio.\\
\textbf{Sustento Técnico:} Sección 3.3 (Impacto en la arquitectura y mantenibilidad del sistema hospitalario).
\begin{tcolorbox}[colback=colorAlex!5,colframe=colorAlex,title=\textbf{Locución Definitiva — Alex}]
Esto transforma por completo nuestro sistema hospitalario. Ahora los módulos de médicos, enfermeros y pacientes están perfectamente organizados en paquetes, documentados a fondo y empaquetados en un único archivo ejecutable. Esto nos permite seguir escalando el sistema médico con total confianza, sabiendo que cada clase cumple con estándares de ingeniería de software limpios y fáciles de mantener.
\end{tcolorbox}

\subsection*{Bloque 11 — 04:29 a 04:42 (13 s) · Referencia: \texttt{11.mp3}}
\textbf{Hablante:} Hanniel \textcolor{colorHanniel}{(Cubo de Geometry Dash, \#5941A9)} \hfill \textbf{Tono:} Cálido, enérgico y cordial\\
\textbf{Visual / Transición:} Gráficos de cierre con el repositorio institucional de GitHub y miembros del equipo despidiéndose.\\
\textbf{Sustento Técnico:} Sección 5 (Conclusiones generales de la práctica).
\begin{tcolorbox}[colback=colorHanniel!5,colframe=colorHanniel,title=\textbf{Locución Definitiva — Hanniel}]
Excelente, Alex. Con paquetes formales, ejecutables portátiles y documentación clara, tenemos una base sólida para seguir creciendo. ¡Muchas gracias por acompañarnos y hasta la próxima entrega!
\end{tcolorbox}

\subsection*{Cierre Visual — 04:42 a 05:00 (18 s)}
\textbf{Visual:} Pantalla final animada con créditos oficiales del Equipo 2, Grupo 7 de POO, enlace al repositorio en GitHub y música de cierre institucional. Sin diálogo hablado.

\end{document}
